% Part: computability
% Chapter: recursive-functions
% Section: pr-functions-computable

\documentclass[../../../include/open-logic-section]{subfiles}

\begin{document}

\olfileid{cmp}{rec}{cmp}
\olsection{আদিম পুনরাবৃত্ত অপেক্ষকগুলি গণনসাধ্য}

ধরা যাক, একটি অপেক্ষক $h$ আদিম পুনরাবৃত্তির সাহায্যে সংজ্ঞায়িত:
\begin{eqnarray*}
h(\vec x, 0)     & = & f(\vec x) \\
h(\vec x, y + 1) & = & g(\vec x, y, h(\vec x, y))
\end{eqnarray*}
এবং ধরা যাক, অপেক্ষক $f$ ও $g$ গণনসাধ্য। ($x_0$, \dots,
$x_{k-1}$-কে সংক্ষেপে বোঝাতে আমরা $\vec x$ ব্যবহার করি।) তাহলে
$h(\vec x, 0)$ অবশ্যই গণনা করা যায়, কারণ সেটি কেবল $f(\vec x)$, যাকে
আমরা গণনসাধ্য বলে ধরে নিয়েছি। তারপর $h(\vec x, 1)$-ও গণনা করা যায়,
কারণ $1 = 0 + 1$ এবং তাই $h(\vec x, 1)$ কেবল
\begin{align*}
 h(\vec x, 1) & = g(\vec x, 0, h(\vec x, 0)) =  g(\vec x, 0, f(\vec x)).
\intertext{এইভাবে এগিয়ে আমরা গণনা করতে পারি}
h(\vec x, 2) & = g(\vec x, 1, h(\vec x, 1)) = g(\vec x, 1, g(\vec x, 0, f(\vec x)))\\
h(\vec x, 3) & = g(\vec x, 2, h(\vec x, 2)) = g(\vec x, 2, g(\vec x, 1, g(\vec x, 0, f(\vec x))))\\
h(\vec x, 4) & = g(\vec x, 3, h(\vec x, 3)) = g(\vec x, 3, g(\vec x, 2, g(\vec x, 1, g(\vec x, 0, f(\vec x)))))\\
& \vdots
\end{align*}
অতএব সাধারণভাবে $h(\vec x, y)$ গণনা করতে আমরা পরপর $h(\vec x, 0)$,
$h(\vec x, 1)$, \dots, গণনা করি, যতক্ষণ না $h(\vec x, y)$-তে পৌঁছাই।

সুতরাং, অপেক্ষক $f$ ও $g$ গণনসাধ্য হলে একটি আদিম পুনরাবৃত্ত সংজ্ঞা
একটি নতুন গণনসাধ্য অপেক্ষক দেয়। একইভাবে, অপেক্ষক $f$ ও $g_i$-গুলি
গণনসাধ্য হলে তাদের মিশ্রণের ফলও একটি গণনসাধ্য অপেক্ষক।

যেহেতু মৌলিক অপেক্ষক $\Zero$, $\Succ$ ও $\Proj{n}{i}$ গণনসাধ্য, এবং
মিশ্রণ ও আদিম পুনরাবৃত্তি গণনসাধ্য অপেক্ষক থেকে গণনসাধ্য অপেক্ষকই দেয়,
তাই প্রত্যেক আদিম পুনরাবৃত্ত অপেক্ষক গণনসাধ্য।

\end{document}
